\usepackage{graphicx} % Required for inserting images
\usepackage[english]{babel}
\usepackage[a4paper]{geometry}
%\usepackage{fancyhdr}
\usepackage[position=top]{subfig}
%\usepackage{multirow}
%\usepackage{lipsum}
%\usepackage{titlesec}
%\usepackage{multicol}
\usepackage{setspace}
%\usepackage{enumitem}
%\usepackage{lipsum}
%\usepackage{rank-2-roots} %!!!!
\usepackage{calc}
\usepackage{enumerate}
\usepackage{centernot}
\usepackage{comment}
%\usepackage{microtype}
%\usepackage{todonotes}
%\usepackage{aeguill} %virgolette francesi
\usepackage{imakeidx}
\usepackage{faktor} %quozienti
\usepackage{stmaryrd} % simboli aggiuntivi


\usepackage{hyperref}
\hypersetup{
    colorlinks,
    citecolor=black,
    %filecolor=black,
    linkcolor=blue,
    urlcolor=black
}

% Teoremi e simboli
\usepackage{amsmath}
\usepackage{yhmath}
\usepackage{mathtools}
\usepackage{amssymb}
\usepackage{amsthm}
\usepackage{mathrsfs}
\usepackage{amsfonts}
%\usepackage{wasysym}
%\usepackage{thmtools}
% extensible arrows
\usepackage{extarrows}

%Box
\newsavebox{\boiteencadrement}
\newenvironment{encadrement}[1][gray!15]
{\par\addvspace{\topsep+10pt}%
    \def\couleurencadrement{#1}%
    \begin{lrbox}{\boiteencadrement}%
        \begin{minipage}{\linewidth-6\fboxsep-1pt}%
            }
            {\end{minipage}\end{lrbox}%
    \noindent\begin{tikzpicture}[
            inner sep=8pt,
            line width=0.75pt]
        \node[rounded corners = 5pt,
            draw,
            color=black,
            fill=\couleurencadrement] {\usebox{\boiteencadrement}};
    \end{tikzpicture}%
    \par\addvspace{\topsep+10pt}%
}

%Equazioni
\numberwithin{equation}{section}

%Teoremi
\theoremstyle{plain}

%\newtheorem{lem}[thm]{Lemma}
%\newtheorem{exer}[thm]{Esercizio}
%\newtheorem{cor}[thm]{Corollary}

\theoremstyle{definition}

\newcounter{thm}
\counterwithin{thm}{section}

%thm
\newtheorem{pretrm}[thm]{Theorem}
\newenvironment{thm}{\begin{encadrement}\begin{pretrm}}{\end{pretrm}\end{encadrement}}

%\newtheorem{defn}[thm]{Definition}
%\newcounter{defn}
%def
\newtheorem{predef}[thm]{Definition}
\newenvironment{defn}{\begin{encadrement}\begin{predef}}{\end{predef}\end{encadrement}}
%prop
\newtheorem{preprop}[thm]{Proposition}
\newenvironment{prop}{\begin{encadrement}\begin{preprop}}{\end{preprop}\end{encadrement}}
%cor
\newtheorem{precor}[thm]{Corollary}
\newenvironment{cor}{\begin{encadrement}\begin{precor}}
            {\end{precor}\end{encadrement}}
%lem
\newtheorem{prelem}[thm]{Lemma}
\newenvironment{lem}{\begin{encadrement}\begin{prelem}}{\end{prelem}\end{encadrement}}
%conj
\newtheorem{preconj}[thm]{Conjecture}
\newenvironment{conj}{\begin{encadrement}\begin{preconj}}{\end{preconj}\end{encadrement}}
%
\newtheorem{exmp}[thm]{Example}
\newtheorem{exer}[thm]{Exercise}
\newtheorem{rmk}[thm]{Remark}
\newtheorem*{rmk*}{Remark}
\newtheorem{facts}[thm]{Facts}
\newtheorem{fact}[thm]{Fact}
\newtheorem{constr}[thm]{Construction}
\newtheorem*{notat}{Notation}
\newtheorem*{note}{Note}


% Colonne
\usepackage{multicol}
\setlength\columnsep{20pt}
\setlength{\columnseprule}{0,5pt}

%% Comandi e operatori
\newcommand{\G}{\Gamma}

% No white line in equal arrows
\usepackage{tikz}
\usetikzlibrary{decorations.markings}
\tikzset{double line with arrow/.style args={#1,#2}{decorate,decoration={markings,%
                    mark=at position 0 with {\coordinate (ta-base-1) at (0,1pt);
                            \coordinate (ta-base-2) at (0,-1pt);},
                    mark=at position 1 with {\draw[#1] (ta-base-1) -- (0,1pt);
                            \draw[#2] (ta-base-2) -- (0,-1pt);
                        }}}}
\tikzset{Equal/.style={-,double line with arrow={-,-}}}

% mathbf
\newcommand{\N}{\mathbf{N}}
\newcommand{\Z}{\mathbf{Z}}
\newcommand{\Q}{\mathbf{Q}}
\newcommand{\QQ}{\overline{\mathbf{Q}}}
\newcommand{\R}{\mathbf{R}}
\newcommand{\K}{K}
\newcommand{\F}{\mathbf{F}}
\newcommand{\C}{\mathbf{C}}
\newcommand{\E}{\mathbf{E}}
% abelianization
\newcommand{\ab}{\text{ab}}


% mathcal
\newcommand{\A}{\mathbf{A}}
\renewcommand{\AA}{\mathcal{A}}
\newcommand{\Sn}{\mathcal{S}_n}
\newcommand{\D}{\mathcal{D}}

\newcommand{\OO}{\mathcal{O}}

% mathfrak
\newcommand{\g}{\mathfrak{g}}
\newcommand{\U}{\mathfrak{U}}
\newcommand{\gl}{\mathfrak{gl}}


\newcommand{\m}{\mathfrak{m}}

%For chapter 3
\newcommand{\cc}{\mathcal{c}}


\renewcommand{\sl}{\mathfrak{sl}}
\renewcommand{\sp}{\mathfrak{sp}}


% math operator
\DeclareMathOperator{\im}{im}
\DeclareMathOperator{\id}{id}
\DeclareMathOperator{\Hom}{Hom}
\DeclareMathOperator{\End}{End}
\DeclareMathOperator{\Aut}{Aut}
\DeclareMathOperator{\Inn}{Inn}
\DeclareMathOperator{\Out}{Out}
\DeclareMathOperator{\coker}{coker}
\DeclareMathOperator{\Stab}{Stab}
\DeclareMathOperator{\Gal}{Gal}
\DeclareMathOperator{\Frob}{Frob}
\DeclareMathOperator{\Span}{Span}
\DeclareMathOperator{\Sym}{Sym}
\DeclareMathOperator{\ad}{ad}
\DeclareMathOperator{\Sp}{Sp}
\DeclareMathOperator{\SO}{SO}
\DeclareMathOperator{\SL}{SL}
\DeclareMathOperator{\GL}{GL}
%\DeclareMathOperator{\OO}{O}
\DeclareMathOperator{\tr}{tr}
%new
\DeclareMathOperator{\car}{char}
\DeclareMathOperator{\Fib}{Fib}
%\DeclareMathOperator{\dep}{dep}
% \legendre{a}{p}
\newcommand{\legendre}[2]{\left(\frac{#1}{#2}\right)}
\renewcommand{\d}{\text{d}}

\newcommand\Iso{\xrightarrow{
        \,\smash{\raisebox{-0.65ex}{\ensuremath{\scriptstyle\sim}}}\,}}
\newcommand\restrict[1]{\raisebox{-.5ex}{$|$}_{#1}}

% Display math
\newcommand{\ssfrac}[2]{
    \raisebox{+0.3ex}{$#1$}
    /
    \raisebox{-0.3ex}{$#2$}
}
% Inline math
\newcommand{\sfrac}[2]{
    \raisebox{+0.3ex}{\scalebox{0.9}{$#1$}}
    /
    \raisebox{-0.3ex}{\scalebox{0.9}{$#2$}}
}

% boldsymbol

\renewcommand{\r}{{\boldsymbol{r}}}
\renewcommand{\k}{{k}}
\newcommand{\h}{{\boldsymbol{h}}}
\newcommand{\x}{{\boldsymbol{x}}}
\newcommand{\0}{{\boldsymbol{0}}}
\newcommand{\s}{{\boldsymbol{s}}}
\renewcommand{\t}{{\boldsymbol{t}}}
\newcommand{\e}{{\boldsymbol{e}}}
%map A_r
%\newcommand{\AR}{{\A}_{\r}}




\newcommand{\tn}{{\otimes n}}
\newcommand{\tens}{\otimes \ldots \otimes}
\newcommand{\n}[1]{\left\lVert #1 \right\lVert}
\newcommand{\res}[1]{\raisebox{-.5ex}{$|$}_{#1}}

%Elliptic curves

\renewcommand{\P}{\mathbf{P}}
\renewcommand{\S}{\operatorname{S}}
\newcommand{\OX}{\OO_{X,P}}
\DeclareMathOperator{\divv}{div}
\DeclareMathOperator{\Pic}{Pic}
\DeclareMathOperator{\Div}{Div}
\renewcommand{\L}{\mathcal{L}}
\newcommand\norm[1]{\left\lVert#1\right\rVert}
\newcommand{\p}{\mathfrak{p}}
\newcommand{\PP}{\mathfrak{P}}
\DeclareMathOperator{\Sel}{Sel}
\newcommand{\Selm}{\Sel^{(m)}}

%Galois groups and fundamental groups
\newcommand{\M}{\mathcal{M}}
% tr deg
\DeclareMathOperator{\trdeg}{tr.deg}
% spec
\DeclareMathOperator{\Spec}{Spec}
% frac
\DeclareMathOperator{\Frac}{Frac}
